% ============================================================
%  SaxoFlow -- SMACD 2026 EDA Competition Paper
%  LaTeX content only (no template preamble).
%  Paste this content block inside your IEEE conference template.
%  Acronym package assumed: \ac{} via acro or acronym package.
%  Bibliography assumed: BibTeX / biblatex with keys listed below.
% ============================================================

% ------------------------------------------------------------------
%  TITLE / AUTHORS  (fill author details before submission)
% ------------------------------------------------------------------
\title{SaxoFlow: An Integrated Open-Source \ac{RTL} Design Suite\\
       with \ac{LLM}-Driven Agentic Automation}

\author{
  \IEEEauthorblockN{First Author\IEEEauthorrefmark{1},
                    Second Author\IEEEauthorrefmark{1}}
  \IEEEauthorblockA{\IEEEauthorrefmark{1}TU Dresden, Dresden, Germany\\
                    \texttt{saxoflowlabs@gmail.com}}
}

\maketitle

% ------------------------------------------------------------------
%  ABSTRACT
% ------------------------------------------------------------------
\begin{abstract}
Open-source \ac{EDA} tool adoption is hindered by fragmented installation
procedures, the absence of a unified project layout, and a steep learning curve
for early-stage designers. SaxoFlow addresses these barriers through a
cohesive platform that integrates fourteen open-source \ac{EDA} tools under a
single \ac{CLI}, a structured project scaffold, and an \ac{LLM}-driven agentic
pipeline capable of generating, reviewing, simulating, and iteratively healing
\ac{RTL} designs and testbenches from a natural-language specification. The
platform pairs a make-based build engine covering simulation, formal
verification, and synthesis with a Rich terminal \ac{UI} featuring AI-assisted
interaction. A provider-agnostic model selector supports eleven \ac{LLM}
backends. The resulting system targets \ac{MSc} and \ac{PhD} students entering
digital design, offering an end-to-end automated path from specification to
verified \ac{RTL}.
\end{abstract}

\begin{IEEEkeywords}
electronic design automation, RTL generation, large language models,
open-source EDA, agentic AI, formal verification, hardware design automation
\end{IEEEkeywords}

% ------------------------------------------------------------------
%  I. INTRODUCTION
% ------------------------------------------------------------------
\section{Introduction}
\label{sec:intro}

The open-source \ac{EDA} ecosystem has matured substantially over the past
decade. Tools such as Yosys~\cite{yosys}, Icarus Verilog~\cite{iverilog},
Verilator~\cite{verilator}, SymbiYosys~\cite{symbiyosys}, and
OpenROAD~\cite{openroad} collectively cover synthesis, simulation, formal
verification, and physical design for both \ac{FPGA} and \ac{ASIC} targets.
Despite this maturity, first-time users face three recurring obstacles. First,
each tool ships with its own installation procedure, version constraints, and
invocation syntax, making a reproducible setup non-trivial. Second, the
absence of a standardised project directory layout forces designers to craft
ad-hoc Makefiles from scratch. Third, translating a design specification into
functional \ac{RTL} requires domain expertise that entry-level students often
lack at the outset.

The rapid adoption of \acp{LLM} in software engineering has prompted parallel
investigations in \ac{HDL} generation~\cite{chipchat,rtlcoder,chateda}.
However, most published work treats code generation as an isolated step, with
no direct feedback from the actual tool chain. A generated Verilog module that
satisfies a grammar checker but fails under Icarus Verilog or Yosys offers
little practical value. The gap between generation and executable verification
remains largely unaddressed.

SaxoFlow closes this gap by coupling an \ac{LLM}-driven multi-agent pipeline
directly with an open-source tool chain. The contributions of the presented
work are as follows:

\begin{itemize}
  \item A unified \ac{CLI} that orchestrates fourteen open-source \ac{EDA}
        tools through preset-based installation profiles and a make-based
        build engine.
  \item A standardised project scaffold that aligns generated artefacts with
        the Makefile targets used for simulation, synthesis, and formal
        verification.
  \item A multi-agent pipeline in which \ac{RTL} and testbench generation,
        structured review, executable simulation, and fault-directed healing
        are chained automatically under a configurable iteration budget.
  \item A provider-agnostic \ac{LLM} selector compatible with eleven backends
        and a Rich terminal \ac{UI} that exposes all capabilities through
        natural-language interaction.
\end{itemize}

The remainder of the paper is structured as follows.
Section~\ref{sec:related} surveys related work.
Section~\ref{sec:methodology} describes the system architecture and
methodology. Section~\ref{sec:conclusion} concludes the paper.

% ------------------------------------------------------------------
%  II. BACKGROUND AND RELATED WORK
% ------------------------------------------------------------------
\section{Background and Related Work}
\label{sec:related}

\subsection{LLM-Assisted Hardware Design}

Several recent efforts apply \acp{LLM} to \ac{HDL} generation.
ChipChat~\cite{chipchat} demonstrated that iterative dialogue with a
conversational model can produce functional Verilog modules, yet the workflow
relied on manual copy-paste cycles between the model and the simulator.
RTL-Coder~\cite{rtlcoder} fine-tuned open-weight models on curated Verilog
corpora, improving first-pass compilation rates, but did not integrate a
downstream verification step. ChatEDA~\cite{chateda} proposed an
\ac{LLM}-assisted flow for \ac{EDA} script generation; however, the scope
was restricted to synthesis scripting rather than full \ac{RTL} design and
verification. AutoChip~\cite{autochip} introduced a feedback loop between a
language model and an \ac{HDL} simulator, which is conceptually closest to
the simulation-grounded healing loop in SaxoFlow, though it did not address
toolchain setup, project scaffolding, or formal verification.

\subsection{Open-Source EDA Frameworks}

At the tool-chain level, projects such as OpenLane~\cite{openlane} provide
end-to-end \ac{ASIC} flows built on OpenROAD and Yosys, but they target
tape-out readiness rather than pedagogical accessibility or \ac{LLM}
integration. No existing framework combines unified installation management,
standardised project layout, \ac{LLM}-driven code generation, and executable
verification in a single coherent system aimed at early-stage digital
designers.

% ------------------------------------------------------------------
%  III. METHODOLOGY
% ------------------------------------------------------------------
\section{System Architecture and Methodology}
\label{sec:methodology}

SaxoFlow comprises three tightly coupled subsystems, illustrated in
Fig.~\ref{fig:architecture}: the \emph{saxoflow} \ac{CLI} and \ac{EDA} flow
engine, the \emph{saxoflow\_agenticai} \ac{LLM} pipeline, and the
\emph{cool\_cli} Rich terminal \ac{UI}. The following subsections describe
each layer in detail.

% ---- Figure 1: Architecture ----
\begin{figure}[t]
  \centering
  % Replace with actual TikZ or exported PDF/PNG diagram
  \includegraphics[width=\columnwidth]{figures/saxoflow_architecture.pdf}
  \caption{SaxoFlow system architecture. The Rich \ac{TUI} delegates commands
           to the unified \ac{CLI}, which drives both the \ac{EDA} toolchain
           and the agentic \ac{LLM} pipeline.}
  \label{fig:architecture}
\end{figure}

\subsection{Unified CLI and EDA Flow Engine}
\label{ssec:cli}

The \texttt{saxoflow} package exposes a Click-based \ac{CLI} whose entry point
is registered as the \texttt{saxoflow} console script. The installer
subsystem groups all supported tools into six named collections -- simulation,
formal, \ac{FPGA}, \ac{ASIC}, base, and \ac{IDE} -- and combines them into
five preset profiles (\texttt{minimal}, \texttt{fpga}, \texttt{asic},
\texttt{formal}, \texttt{full}). Each profile is resolved at install time to
either an \texttt{apt-get} package or a shell recipe located under
\texttt{scripts/recipes/}. Script-installed binaries are appended to the
virtual environment activation script, ensuring correct \texttt{PATH}
resolution on subsequent invocations.

Tool health is monitored through a diagnostics subsystem that infers the
active flow profile from the saved tool selection, probes \texttt{PATH} and
common \texttt{\$HOME/.local} prefixes, extracts version strings via
tool-specific regular expressions, and computes a weighted health score
normalised over required and optional tool sets.
Table~\ref{tab:tools} summarises the supported tools across categories.

The \texttt{saxoflow unit} command scaffolds a deterministic project tree
containing separate directories for \ac{RTL} sources (Verilog, \ac{SV}, VHDL),
testbenches, simulation outputs, synthesis scripts and reports, formal
verification scripts and reports, and constraint files. A pre-filled Yosys
synthesis script and a universal Makefile are placed at the project root. This
layout is shared between the manual flow and the agentic pipeline, ensuring
that automatically generated artefacts land in paths that Makefile targets
already reference.

The build engine (\texttt{saxoflow/makeflow.py}) provides Click commands for
Icarus Verilog simulation, Verilator \texttt{C++} build and execution, GTKWave
waveform viewing, SymbiYosys formal verification, Yosys synthesis, and
housekeeping. Testbench resolution is automated: if a single testbench file is
present, it is selected without user interaction; if multiple files exist,
a numbered prompt is displayed.

% ---- Table I: Tools ----
\begin{table}[t]
  \caption{Supported EDA Tool Integration in SaxoFlow}
  \label{tab:tools}
  \centering
  \scriptsize
  \begin{tabular}{llll}
    \hline
    \textbf{Category} & \textbf{Tool} & \textbf{Install} & \textbf{Purpose} \\
    \hline
    Simulation  & iverilog          & APT    & Verilog-2001 event-driven sim \\
    Simulation  & Verilator         & Script & Cycle-accurate SV-to-C++ sim \\
    Waveform    & GTKWave           & APT    & VCD/FST waveform viewer \\
    Synthesis   & Yosys+slang       & Script & RTL-to-gate synthesis \\
    Formal      & SymbiYosys        & Script & Formal property checking \\
    FPGA PnR    & nextpnr           & Script & Open-source place-and-route \\
    FPGA prog.  & openFPGAloader    & APT    & FPGA bitstream uploader \\
    FPGA vendor & Vivado            & Script & Xilinx FPGA suite (optional) \\
    ASIC PD     & OpenROAD          & Script & Digital ASIC backend flow \\
    ASIC layout & KLayout           & APT    & GDS viewer and scripting \\
    ASIC layout & Magic             & APT    & VLSI layout editor \\
    ASIC LVS    & Netgen            & APT    & Layout-vs-schematic check \\
    HDL deps    & Bender            & Script & HDL dependency manager \\
    IDE         & VS Code           & Script & RTL development environment \\
    \hline
  \end{tabular}
\end{table}

\subsection{LLM-Driven Agentic Pipeline}
\label{ssec:agentic}

The \texttt{saxoflow\_agenticai} module implements a multi-agent pipeline
centred on a generate-review-improve loop. Nine agents are registered in a
factory (\texttt{AgentManager}): four generators (RTL, testbench, formal
properties, report), four reviewers (RTL review, testbench review, formal
property review, debug), and one \ac{EDA}-tool agent (simulation). The
complete automated flow is orchestrated by \texttt{AgentOrchestrator.full\_pipeline()},
whose phases are shown in Fig.~\ref{fig:pipeline}.

% ---- Figure 2: Agentic pipeline ----
\begin{figure}[t]
  \centering
  % Replace with actual TikZ flowchart
  \includegraphics[width=\columnwidth]{figures/saxoflow_pipeline.pdf}
  \caption{Agentic \ac{LLM} pipeline. Each generation phase is followed by a
           structured review loop. Simulation failure triggers the debug agent,
           which selects a healing agent for the next iteration.}
  \label{fig:pipeline}
\end{figure}

\textbf{Abstract base agent.}
All \ac{LLM}-driven agents extend \texttt{BaseAgent}, an abstract base class
that handles prompt template loading (cached after the first access),
Jinja2-based template rendering via LangChain's \texttt{PromptTemplate} with
\texttt{template\_format="jinja2"}, and model invocation through
\texttt{llm.invoke()}. Structured critique and code extraction are handled by
subclass-specific post-processing functions using regular-expression stripping
of markdown fence artefacts.

\textbf{Feedback coordinator.}
\texttt{AgentFeedbackCoordinator.iterate\_improvements()} implements the
generate-review-improve loop. After an initial generation pass, the
corresponding review agent produces a critique. An early-exit predicate
(\texttt{is\_no\_action\_feedback}) matches eleven approval phrases (e.g.,
\emph{no issues}, \emph{looks good}, \emph{approved}) to terminate the loop
before the iteration budget is exhausted. Otherwise, the generator's
\texttt{improve()} method is invoked with the review as additional context.

\textbf{Simulation-grounded healing.}
The simulation agent wraps \texttt{saxoflow.makeflow.sim} via Click's
\texttt{CliRunner} inside a directory-change context manager, capturing
stdout and stderr and checking for the presence of a \ac{VCD} output to
determine success. On failure, the debug agent analyses the compiler and
simulation error messages and returns a structured recommendation identifying
which generator agents (\texttt{RTLGenAgent}, \texttt{TBGenAgent}, or
\texttt{UserAction}) should be applied in the next healing iteration. The
heal-and-resimulate cycle repeats up to a configurable maximum iteration count.
This closes the feedback loop between the \ac{LLM} and the actual tool chain,
preventing the generation of code that is syntactically plausible but
functionally or tool-incompatibly incorrect.

\textbf{Prompt engineering strategy.}
A two-layer prompt architecture is employed. The first layer prepends
tool-constraint guidelines to every prompt: \texttt{verilog\_guidelines.txt}
and \texttt{verilog\_constructs.txt} encode the allowed and forbidden
Verilog-2001 construct set for Icarus, Verilator, Yosys, and OpenROAD
compatibility; \texttt{tb\_guidelines.txt} and \texttt{tb\_constructs.txt}
encode analogous rules for testbench generation. The second layer contains
the task-specific Jinja2 template, parameterised with \texttt{\{\{spec\}\}},
\texttt{\{\{rtl\_code\}\}}, and conditional \texttt{\{\{review\}\}} blocks.
All prompts instruct the model to output only a bare code block bounded by
\texttt{module}/\texttt{endmodule}, with no markdown fencing or explanatory
text, enabling deterministic post-processing.

\textbf{Provider-agnostic model selection.}
\texttt{ModelSelector} supports eleven \ac{LLM} providers spanning OpenAI,
Anthropic, Google Gemini, Groq, Mistral, Fireworks, Together, Perplexity,
DeepSeek, DashScope, and OpenRouter. At startup, the selector scans an
autodetect priority list for the first provider whose \ac{API} key is set
in the environment and constructs the appropriate LangChain adapter:
\texttt{ChatOpenAI} for OpenAI-compatible endpoints (with custom
\texttt{base\_url} and headers where required), \texttt{ChatAnthropic} for
Anthropic, and \texttt{ChatGoogleGenerativeAI} for Gemini. A \ac{YAML}
configuration file allows per-agent model and temperature overrides. This
design decouples the pipeline from any single provider, reducing the barrier
to adoption across different institutional or personal \ac{API} access
scenarios.

\subsection{Rich Terminal User Interface}
\label{ssec:tui}

The \texttt{cool\_cli} package provides a Rich-based terminal \ac{UI} launched
via \texttt{python3 start.py}. An interactive prompt loop built on
\texttt{prompt\_toolkit} accepts free-form input and routes it to one of four
handlers: built-in commands (\texttt{help}, \texttt{clear}, \texttt{quit}),
agentic commands dispatched as subprocesses to the \texttt{saxoflow\_agenticai}
\ac{CLI}, Unix shell aliases (\texttt{ls}, \texttt{pwd}, editor invocations),
and the AI Buddy, a conversational assistant backed by the same
\texttt{ModelSelector} infrastructure.

The AI Buddy maintains a sliding window of five conversation turns and
performs keyword-based intent detection over forty action mappings
(e.g., \emph{generate rtl}, \emph{review testbench}, \emph{simulate}).
On detection of a design-related intent, it requests the relevant source file
from the user, invokes the appropriate review or generation agent, and returns
a structured panel-formatted response. A bootstrap routine runs at first
launch, creating a \texttt{.env} template and guiding the user interactively
through \ac{API} key configuration for the detected provider.

% ------------------------------------------------------------------
%  IV. CONCLUSION
% ------------------------------------------------------------------
\section{Conclusion}
\label{sec:conclusion}

SaxoFlow presents an integrated solution to the fragmentation of open-source
\ac{EDA} tool usage by combining a unified installation and build framework
with a simulation-grounded \ac{LLM} pipeline and an accessible terminal
\ac{UI}. The key distinguishing characteristic of the approach is the closure
of the feedback loop between \ac{LLM}-generated \ac{HDL} and the actual
open-source tool chain: generated \ac{RTL} and testbenches are compiled and
simulated, simulator diagnostics are analysed by a dedicated debug agent, and
targeted healing is applied iteratively. The two-layer prompt architecture
enforces tool-specific construct constraints at the language model level,
reducing tool-incompatible outputs prior to any compilation attempt.

Future directions include quantitative benchmarking across a curated set of
design specifications, re-enabling the formal property generation phase within
the automated pipeline, extending the installer to non-Debian distributions,
and investigating parallel agent execution to reduce end-to-end pipeline
latency.

% ------------------------------------------------------------------
%  REFERENCES
%  Replace placeholder keys with actual BibTeX entries.
% ------------------------------------------------------------------
\begin{thebibliography}{10}

\bibitem{yosys}
C.~Wolf, \emph{Yosys Open Synthesis Suite}, 2013, [Online].
Available: \url{https://yosyshq.net/yosys/}

\bibitem{iverilog}
S.~Williams, \emph{Icarus Verilog}, 2002, [Online].
Available: \url{http://iverilog.icarus.com/}

\bibitem{verilator}
W.~Snyder, \emph{Verilator: The Fast Free Verilog/SystemVerilog Simulator},
in \emph{Proc. DesignCon}, 2013.

\bibitem{symbiyosys}
C.~Wolf, \emph{SymbiYosys: Formal Hardware Verification Front-End},
2017, [Online].
Available: \url{https://symbiyosys.readthedocs.io/}

\bibitem{openroad}
T.~Ajayi \emph{et al.}, ``OpenROAD: Toward a Self-Driving,
Open-Source Digital Layout Implementation Tool Chain,''
in \emph{Proc. GOMACTECH}, 2019, pp. 1105--1110.

\bibitem{chipchat}
J.~Blocklove \emph{et al.}, ``ChipChat: Have LLMs Become Good Hardware
Designers?'' in \emph{Proc. ICCAD}, 2023.

\bibitem{rtlcoder}
S.~Liu \emph{et al.}, ``RTLCoder: Outperforming GPT-3.5 in Design RTL
Generation with Our Open-Source Dataset and Lightweight Solution,''
\emph{arXiv:2312.08617}, 2023.

\bibitem{chateda}
Z.~He \emph{et al.}, ``ChatEDA: A Large Language Model Powered Autonomous
Agent for EDA,'' \emph{IEEE Trans.\ Comput.-Aided Design Integr.\ Circuits
Syst.}, vol.~43, no.~9, pp. 2946--2959, 2024.

\bibitem{autochip}
B.~Tsai \emph{et al.}, ``AutoChip: Automating HDL Generation Using
LLM Feedback,'' \emph{arXiv:2311.04887}, 2023.

\bibitem{openlane}
M.~Shalan and T.~Edwards, ``Building OpenLANE: A 130nm OpenROAD-Based
Tapeout-Proven Flow,'' in \emph{Proc. ICCAD}, 2020.

\end{thebibliography}
